\section{总结与延伸}

\subsection{五种实现方式的权衡}

\begin{table}[H]
\centering
\begin{tabular}{lccccc}
\toprule
& \textbf{Manual} & \textbf{PyTorch} & \textbf{torch.compile} & \textbf{CUDA} & \textbf{Triton} \\
\midrule
开发难度 & 最低 & 低 & 低 & 高 & 中 \\
性能 & 最差 & 好 & 很好 & 好--很好 & 好--很好 \\
灵活性 & 高 & 受限 & 中等 & 最高 & 高 \\
可调试性 & 好 & 一般 & 差 & 差 & 好 \\
自动融合 & 否 & 已融合 & 是 & 否 & 否 \\
\bottomrule
\end{tabular}
\caption{五种实现方式的多维度对比}
\end{table}

\subsection{全课知识图谱}

\begin{center}
\begin{tikzpicture}[
    node distance=0.9cm,
    every node/.style={draw, rounded corners, minimum width=3.2cm, minimum height=0.7cm, font=\small}
]
\node (bench) {Benchmarking \& Profiling};
\node (fusion) [below=of bench] {Kernel Fusion 动机};
\node (cuda) [below left=1cm and -0.5cm of fusion] {CUDA 内核};
\node (triton) [below=of fusion] {Triton 内核};
\node (compile) [below right=1cm and -0.5cm of fusion] {torch.compile};
\node (softmax) [below=of triton] {Softmax（行归约）};
\draw[->, thick] (bench) -- (fusion);
\draw[->, thick] (fusion) -- (cuda);
\draw[->, thick] (fusion) -- (triton);
\draw[->, thick] (fusion) -- (compile);
\draw[->, thick] (triton) -- (softmax);

\node[draw=none, right=1.2cm of bench, text width=4.5cm, font=\scriptsize] {warmup, synchronize, profiler, flame graph};
\node[draw=none, right=1.2cm of fusion, text width=4.5cm, font=\scriptsize] {仓库-工厂类比, GeLU 8x 加速};
\node[draw=none, left=1.2cm of softmax, text width=4.5cm, font=\scriptsize] {一行=一Block, mask 处理边界};
\end{tikzpicture}
\end{center}

\subsection{关键 Takeaways}

\begin{importantbox}{六条核心原则}
\begin{enumerate}
    \item \textbf{编程模型与硬件之间有巨大鸿沟}：性能之谜源于此，理解这一点是性能优化的前提
    \item \textbf{始终基准测试和 Profile}：不要凭直觉判断瓶颈，让数据说话
    \item \textbf{核心原则是最小化读写}：内核融合、tiling、重计算本质上都是在减少全局内存访问
    \item \textbf{torch.compile 是低垂的果实}：对标准操作，优先使用编译器自动优化
    \item \textbf{Triton 是自定义内核的首选}：Python 开发体验 + 接近 CUDA 的性能
    \item \textbf{自动编译器会越来越好}：今天需要手写的优化，明天可能由编译器自动完成
\end{enumerate}
\end{importantbox}

\subsection{拓展阅读}

\begin{itemize}
    \item Horace He 博客: \url{https://horace.io/brrr_intro.html} —— GPU 性能优化入门必读
    \item Triton 官方教程（Fused Softmax）: \url{https://triton-lang.org/main/getting-started/tutorials/02-fused-softmax.html}
    \item Triton 官方教程（Matrix Multiplication）: \url{https://triton-lang.org/main/getting-started/tutorials/03-matrix-multiplication.html}
    \item Triton 论文: Tillet et al., \url{https://www.eecs.harvard.edu/~htk/publication/2019-mapl-tillet-kung-cox.pdf}
    \item CUDA MODE Lecture 1（如何 Profile CUDA 内核）: \url{https://www.youtube.com/watch?v=LuhJEEJQgUM}
    \item CUDA MODE Lecture 2（PPMP 书 ch.1--3）: \url{https://www.youtube.com/watch?v=NQ-0D5Ti2dc}
    \item CUDA MODE Lecture 3（CUDA for Python Programmers）: \url{https://www.youtube.com/watch?v=4sgKnKbR-WE}
    \item CUDA MODE Lecture 4（Compute and Memory Basics）: \url{https://www.youtube.com/watch?v=lTmYrKwjSOU}
    \item CUDA MODE Lecture 8（CUDA Performance Checklist）: \url{https://www.youtube.com/watch?v=SGhfUhlowB4}
    \item HetSys Course（GPU Programming）: \url{https://www.youtube.com/watch?v=8JGo2zylE80}
    \item GPU Puzzles: \url{https://github.com/srush/gpu-puzzles}
    \item NVIDIA Deep Learning Performance Guide: \url{https://docs.nvidia.com/deeplearning/performance/dl-performance-gpu-background/index.html}
    \item PyTorch 2.0 Operator Fusion:\\ {\small\url{https://towardsdatascience.com/how-pytorch-2-0-accelerates-deep-learning-with-operator-fusion-and-cpu-gpu-code-generation-35132a85bd26}}
\end{itemize}

\end{document}
